% Deck for trying pdf-present: speaker notes on the right half of each
% page.  Built by ./build.sh.
\documentclass[aspectratio=169]{beamer}
\setbeamertemplate{navigation symbols}{}
\setbeameroption{show notes on second screen=right}
\title{Presenting from Emacs}
\subtitle{A demo deck for \texttt{pdf-present}}
\date{}
\begin{document}

\begin{frame}
  \titlepage
  \note{Welcome. The presenter window shows this note, the next slide, the elapsed time and the clock. Press d to set a talk length and watch the countdown.}
\end{frame}

\begin{frame}{Two windows, not two fullscreen frames}
  \begin{itemize}
  \item The audience window shows the slide alone
  \item The presenter window shows this slide, the next one and the notes
  \item Both are plain windows, so Wayland compositors place them
  \end{itemize}
  \note{Point out that nothing re-renders when the windows settle: each is made in its final size.}
\end{frame}

\begin{frame}{Overlays count as slides}
  \begin{itemize}
  \item<1-> First step
  \item<2-> Second step, on the next page of the PDF
  \item<3-> Third step
  \end{itemize}
  \note{Press n three times here. The slide counter uses the page labels, so it says 3 for all three steps.}
\end{frame}

\begin{frame}{Jumping around}
  \begin{itemize}
  \item \texttt{g} goes to a slide by number
  \item \texttt{T} shows a thumbnail grid; the audience keeps the current slide
  \item \texttt{<} and \texttt{>} go to the first and last slides
  \end{itemize}
  \note{Try T now and pick the title slide, then come back with g.}
\end{frame}

\begin{frame}{Blank the screen}
  \centering\Large Press \texttt{b} for black, \texttt{w} for white
  \note{Useful when the class should look at the board, not the slide. Any key brings the slide back.}
\end{frame}

\begin{frame}{The deck reloads when it changes}
  Edit this slide in \texttt{demo-slides.tex}, run \texttt{./build.sh}, and the slides re-read in place, keeping the current slide and the timers.
  \note{This is the edit-while-presenting loop: fix a typo during questions without restarting the talk.}
\end{frame}

\begin{frame}
  \centering\Huge Questions?
  \note{q quits and restores your buffer, windows and frames.}
\end{frame}

\end{document}
